\documentclass[../../../include/open-logic-section]{subfiles}

\begin{document}

\olfileid{his}{set}{cantorplane}
\olsection{كانتور والخط والمستوى}

لبعض ما أحاط ببرهان شرودر--برنشتاين صلة بالتاريخ المتقدم في
~\olref[his][set][pathology]{sec}. فقد أثبت كانتور سنة 1877 أن نقاط
المربع مساوية في العدد لنقاط ضلع واحد منه. وسنورد ها هنا برهانه، أو قل
محاولته الأولى للبرهان.

لنسمّ~$\unitline$ قطعة الوحدة المستقيمة، أي مجموعة النقاط~$[0,1]$،
ولنسمّ~$\unitsquare$ مربع الوحدة، أي مجموعة النقاط~$\unitline \times
\unitline$. فقول كانتور بهذه العبارة هو
$\cardeq{\unitline}{\unitsquare}$. وقد أرسل إلى ديدكند رسالة حاصل حجتها
ما يأتي.

\begin{thm}\ollabel{thm:cantorplane}
$\cardeq{\unitline}{\unitsquare}$
\end{thm}

\begin{proof}[البرهان: الجزء الأول.]
خذ~$a, b \in \unitline$، واكتبهما بالترميز الثنائي؛ فتحصل متتاليتان لا
نهاية لهما من الأصفار والآحاد،~$a_1$،~$a_2$،~\dots، و$b_1$،~$b_2$،
~\dots، على الوجه الآتي:
\begin{align*}
a &= 0.a_1a_2a_3a_4\dots\\
b &= 0.b_1b_2b_3b_4\dots
\intertext{تأمل الآن الدالة~$f \colon \unitsquare \to \unitline$ المعطاة بـ}
f(a, b) & = 0.a_1b_1a_2b_2a_3b_3a_4b_4\dots
\end{align*}
فتكون~$f$ !!a{injection}؛ لأن من~$f(a, b) = f(c,d)$ يلزم
$a_n = c_n$ و$b_n = d_n$ لكل~$n \in \Nat$، فيلزم~$a = c$ و$b = d$.
\end{proof}

غير أن ديدكند بيّن لكانتور أن هذا لا يفي بالسؤال الأول. فانظر إلى
$0.\dot{1}\dot{0} = 0.1010101010\ldots$. ينبغي أن نجد
$f(a,b) = 0.\dot{1}\dot{0}$، مع:
\begin{align*}
a&= 0.\dot{1}\dot{1} = 0.111111\ldots\\
b&= 0
\end{align*}
ولكن~$a = 0.\dot{1}\dot{1} = 1$. فإذا قلنا «اكتب~$a$ و$b$ بالترميز
الثنائي»، فلا بد من تعيين \emph{أي} التمثيلين نأخذ؛ ولأن~$f$ لا بد أن
تكون \emph{دالة}، وجب أن نلتزم تمثيلاً \emph{واحداً} من التمثيلين
الجائزين. فإذا التزمنا، مثلاً، التمثيل المنتهي وكتبنا~$a$ هكذا
«$1.000\ldots$»، لم يوجد زوج~$\tuple{a, b}$ يحقق
$f(a, b) = 0.\dot{1}\dot{0}$.

وحاصل تنبيه ديدكند أن وجود بعض التوسعات الثنائية الدورية يجعل دالة
كانتور~$f$ !!a{injection}، ولكنها \emph{ليست} !!a{surjection}. فلم يثبت
كانتور بهذا إلا~$\cardle{\unitsquare}{\unitline}$، دون
$\cardeq{\unitsquare}{\unitline}$.

فأجاب كانتور ديدكند من فوره أو كاد، وقال في حاصل كلامه إن البرهان يتم على
هذا الوجه:

\begin{proof}[البرهان: الإتمام.]
قد ثبت~$\cardle{\unitsquare}{\unitline}$. ومن البيّن وجود !!a{injection}
من~$\unitline$ إلى~$\unitsquare$؛ فما عليك إلا أن تبسط القطعة المستقيمة
على ضلع من أضلاع المربع. فيكون
$\cardle{\unitline}{\unitsquare}$ و
$\cardle{\unitsquare}{\unitline}$. ثم بمبرهنة شرودر--برنشتاين
(\olref[sfr][siz][sb]{thm:schroder-bernstein})، نحصل على
$\cardeq{\unitline}{\unitsquare}$.
\end{proof}

ولم يكن في وسع كانتور، بداهة، أن يختم كلامه يومئذ بهذه العبارة، إذ لم تكن
مبرهنة شرودر--برنشتاين قد ثبتت بعد. بل إنه، وإن صاغها فيما بعد حدساً
عاماً، لم يقف على برهان يرضاه لها حتى سنة 1897. (ولهذا أتى كانتور، في
وقت لاحق من سنة 1877، ببرهان آخر لـ\olref{thm:cantorplane} لا يمر
بشرودر--برنشتاين.)

\end{document}
